
\bl{El siguiente resultado es un ejercicio en el libro \cite[Exer 1.1-1.5]{kuipers-1974}, cuya incluimos aqu\'i por completitud. }

\begin{proposition}
    \label{def:equivalentes}
    Sea $x=(x_n)_{n\geq1}$ una sucesión de números reales.
    Las siguientes afirmaciones son equivalentes:
    \begin{enumerate}
        \item $(x_n)$ es equidistribuida módulo $1$.

        \item La igualdad (\ref{def:equidistribucion}) se cumple
        para todo intervalo $(a,b)$ con $0 \leq a < b \leq 1$.

        \item La igualdad (\ref{def:equidistribucion}) se cumple
        para todo intervalo $[a,b], \,[a,b), \,(a,b]$ con $0 \leq a < b \leq 1$.

        \item La igualdad (\ref{def:equidistribucion}) se cumple
        para todo intervalo $[0,b)$ con $0 \leq b \leq 1$.

        \item La igualdad (\ref{def:equidistribucion}) se cumple
        para todo intervalo $(a,1)$ con $0 \leq a \leq 1$.

        \item La igualdad (\ref{def:equidistribucion}) se cumple
        para todo intervalo $(a,b)$ con $a,b\in D$, donde $D \subseteq \mathbb{T}$ es un conjunto denso.
    \end{enumerate}
\end{proposition}

\begin{proof}
    $1 \Longrightarrow 2$ \, Supongamos que $(x_n)$ es equidistribuida módulo $1$,
    es decir, que (\ref{def:equidistribucion}) se cumple para todo intervalo
    $(a,b) \subset \mathbb{T}$ con $0 < a < b < 1$.

    Basta probar que también se cumple para intervalos que contienen los extremos.

    \textit{Caso $(0,b)$}: \, Sea $\varepsilon > 0$, entonces
    \[
        A_N((\varepsilon, b), x) \leq A_N((0,b), x) 
        \quad \text{y} \quad 
        A_N((0,b), x) + A_N((b, 1-\varepsilon), x) \leq N
    \]
    Dividiendo por $N$ y tomando el límite cuando $N \to \infty$, 
    \[
        b-\varepsilon \leq 
        \liminf_{N \to \infty} \frac{A_N((0,b), x)}{N} \leq
        \limsup_{N \to \infty} \frac{A_N((0,b), x)}{N} \leq 
        1 - (1-b-\varepsilon) = 
        b + \varepsilon
    \]
    Como $\varepsilon > 0$ es arbitrario, se deduce que el límite existe y es igual a $b$.
    
    \textit{Caso $(a,1)$}: \, Sea $\varepsilon > 0$, entonces
    \[
        A_N((a, 1-\varepsilon), x) \leq A_N((a,1), x) 
        \quad \text{y} \quad 
        A_N((\varepsilon,a), x) + A_N((a, 1), x) \leq N
    \]
    Dividiendo por $N$ y tomando el límite cuando $N \to \infty$, 
    \[
        1-a-\varepsilon \leq 
        \liminf_{N \to \infty} \frac{A_N((a,1), x)}{N} \leq
        \limsup_{N \to \infty} \frac{A_N((a,1), x)}{N} \leq 
        1 - (a-\varepsilon) = 
        1- a + \varepsilon
    \]
    Como $\varepsilon > 0$ es arbitrario, se deduce que el límite existe y es igual a $1-a$.

    $2 \Longrightarrow 1$ \, Es obvio.

    $2 \Longrightarrow 3$ \, Veamos que para $0\leq c \leq 1$ se tiene que
    \[
        \lim_{N \to \infty} \frac{\card \left\{1\leq n \leq N :\partfrac{x_n}=c\right\}}{N} = 0
    \]
    Sea $\varepsilon > 0$, si $c \in (0,1)$, entonces
    \[
        0 \leq 
        \limsup_{N \to \infty} \frac{\card \left\{1\leq n \leq N :\partfrac{x_n}=c\right\}}{N} \leq
        \lim_{N \to \infty} \frac{A_N((c-\varepsilon, c+\varepsilon), x)}{N} = 2\varepsilon
    \]
    como $\varepsilon > 0$ es arbitrario, el límite existe y es igual a $0$.
    
    Sea $\varepsilon > 0$, si $c = 0$, entonces
    \[
        0 \leq 
        \card \left\{1\leq n \leq N :\partfrac{x_n}=0\right\} +
        A_N((\varepsilon, 1), x) \leq
        N
    \]
    luego
    \[
        0 \leq 
        \limsup_{N \to \infty} \frac{\card \left\{1\leq n \leq N :\partfrac{x_n}=0\right\}}{N} \leq
        1 - (1-\varepsilon) = 
        \varepsilon
    \]
    como $\varepsilon > 0$ es arbitrario, el límite existe y es igual a $0$.

    Sea $\varepsilon > 0$, si $c = 1$, entonces
    \[
        0 \leq 
        \card \left\{1\leq n \leq N :\partfrac{x_n}=1\right\} +
        A_N((0, 1-\varepsilon), x) \leq
        N
    \]
    luego
    \[
        0 \leq 
        \limsup_{N \to \infty} \frac{\card \left\{1\leq n \leq N :\partfrac{x_n}=1\right\}}{N} \leq
        1 - (1-\varepsilon) = 
        \varepsilon
    \]  
    como $\varepsilon > 0$ es arbitrario, el límite existe y es igual a $0$.

    Por lo tanto, para los conjuntos unipuntuales tenemos que su limite es $0$.
    Como los intervalos $[a,b], \,[a,b), \,(a,b]$ se descomponen de la forma
    \[
        [a,b] = (a,b) \cup \{a\} \cup \{b\} \quad
        [a,b) = (a,b) \cup \{a\} \quad
        (a,b] = (a,b) \cup \{b\}
    \]
    entonces la igualdad (\ref{def:equidistribucion}) se cumple para estos intervalos.
    
    $3 \Longrightarrow 2$ \, Usando el mismo argumento para cada uno de los tipos de intervalos
    se puede demostrar que para $0 \leq c \leq 1$,
    \[
        \lim_{N \to \infty} \frac{\card \left\{1\leq n \leq N :\partfrac{x_n}=c\right\}}{N} = 0
    \]
    y teniendo en cuenta la descomposición se obtiene la igualdad para intervalos de la forma $(a,b)$.

    De hecho, los diferentes intervalos $[a,b], \,[a,b), \,(a,b]$ son equivalentes entre sí, 
    es decir, que si la igualdad (\ref{def:equidistribucion}) se cumple para un tipo de 
    intervalo entonces se cumple para los otros tipos.
    
    $3 \Longrightarrow 4$ \, Es obvio.

    $4 \Longrightarrow 3$ \, Sea $[a,b)$ fijo, como $[a,b) = [0,b) \setminus [0,a)$ entonces 
    \[
        \lim_{N \to \infty} \frac{A_N([a,b), x)}{N} = 
        \lim_{N \to \infty} \left(\frac{A_N([0,b), x)}{N} - \frac{A_N([0,a), x)}{N}\right) = 
        b - a
    \]
    como es cierto para todo $[a,b)$, entonces también es cierto para todo $[a,b], \,(a,b]$.

    $2 \Longrightarrow 5$ \, Es obvio.

    $5 \Longrightarrow 3$ \, Sea $(a,b]$ fijo, como $(a,b]=(a,1) \setminus (b,1)$, entonces
    \[
        \lim_{N \to \infty} \frac{A_N((a,b], x)}{N} = 
        \lim_{N \to \infty} \left(\frac{A_N((a,1), x)}{N} - \frac{A_N((b,1), x)}{N}\right) = 
        b - a
    \]
    como es cierto para todo $(a,b]$, entonces también es cierto para todo $[a,b], \,[a,b)$.
    
    $2 \Longrightarrow 6$ \, Es obvio.

    $6 \Longrightarrow 1$ \, Sea $(a,b)$ fijo y sea $\varepsilon > 0$. Por densidad de $D$,
    existen $\alpha, \beta \in D$ tales que $a-\varepsilon < \alpha < a$ y $b < \beta < b+\varepsilon$.
    Entonces
    \[
        \limsup_{N \to \infty} \frac{A_N((a,b), x)}{N} \leq
        \lim_{N \to \infty} \frac{A_N((\alpha,\beta), x)}{N} = 
        \beta - \alpha <
        b - a + 2\varepsilon
    \]
    como $\varepsilon > 0$ es arbitrario, se deduce que 
    \[
        \limsup_{N \to \infty} \frac{A_N((a,b), x)}{N} \leq b - a
    \]
    De forma análoga, se obtiene que 
    \[
        \liminf_{N \to \infty} \frac{A_N((a,b), x)}{N} \geq b - a
    \]
    En consecuencia, 
    \[
        \lim_{N \to \infty} \frac{A_N((a,b), x)}{N} = b - a
    \]

    Por lo tanto, todas las afirmaciones son equivalentes entre sí.
\end{proof}